\section{Marco teórico}

El sistema desarrollado en este proyecto implementa un monitor de ECG en tiempo real basado en cuatro etapas principales: adquisición de datos, preprocesamiento digital, detección de picos R y visualización interactiva. El objetivo es obtener una representación clara de la actividad eléctrica cardiaca, estimar la frecuencia cardiaca a partir de los intervalos entre latidos y mostrar simultáneamente la evolución de la señal en sus distintas etapas de filtrado. La solución combina principios de instrumentación biomédica, tratamiento digital de señales y programación orientada a interfaces gráficas.

\subsection{Señal electrocardiográfica}

El electrocardiograma es el registro de la actividad eléctrica del corazón medido en la superficie corporal. La señal ECG refleja la propagación de la despolarización y repolarización cardiaca, y suele describirse mediante las ondas P, complejo QRS y onda T. Dentro de estas componentes, el complejo QRS es el más relevante para el monitoreo de ritmo, debido a que posee la mayor amplitud y marca la despolarización ventricular. En particular, el pico R, que corresponde al máximo principal del complejo QRS, se emplea como referencia para contar latidos y calcular la frecuencia cardiaca.

Desde el punto de vista fisiológico, una señal ECG útil debe preservar la morfología del QRS, ya que pequeñas variaciones en su forma, amplitud o separación temporal pueden indicar cambios en el estado del paciente. Sin embargo, la señal adquirida en condiciones reales suele verse afectada por ruido muscular, desplazamiento de línea base, interferencia de la red eléctrica, artefactos de movimiento y variaciones debidas al contacto de los electrodos. Por ello, antes de analizarla es necesario aplicar filtrado y criterios de detección robustos.

\subsection{Adquisición y flujo de datos en tiempo real}

La adquisición en este proyecto se realiza desde un microcontrolador que envía muestras discretas por comunicación serie. Cada valor recibido representa una lectura instantánea de la señal ECG y se procesa de forma incremental, es decir, a medida que llega. Esta estrategia es propia de los sistemas de monitoreo en vivo, donde no se dispone de toda la señal completa sino de fragmentos sucesivos que deben integrarse al instante.

Para mantener la interfaz gráfica fluida, la lectura del puerto serie se ejecuta en un hilo independiente al de la ventana principal. De este modo, la recepción de muestras no bloquea el repintado de las gráficas ni la actualización del indicador de frecuencia cardiaca. En paralelo, el sistema también contempla un modo de simulación mediante archivos CSV previamente capturados, lo cual permite reproducir la señal como si proviniera del hardware real.

\subsection{Preprocesamiento y filtrado causal}

El tratamiento de la señal se fundamenta en un filtrado causal, es decir, un filtrado que utiliza únicamente la muestra actual y las anteriores, sin depender de datos futuros. Esta propiedad es imprescindible en un sistema en tiempo real, porque permite procesar la señal de manera continua conforme se reciben nuevos datos. A diferencia de los métodos no causales, el filtrado causal es compatible con la visualización en vivo y con la estimación instantánea de eventos cardiacos.

En el sistema se emplea una cadena de filtrado de dos etapas. La primera corresponde a un filtro pasa-banda de 0.5 a 40 Hz, cuyo propósito es conservar el contenido principal de interés fisiológico y atenuar componentes no deseadas. El límite inferior reduce la deriva de la línea base, que suele aparecer por respiración o cambios en el contacto de los electrodos, mientras que el límite superior atenúa ruido de alta frecuencia asociado a actividad muscular o interferencias de adquisición. En el contexto del ECG, este rango es adecuado porque conserva la energía principal del complejo QRS y mejora la legibilidad general de la señal.

La segunda etapa es un filtro \textit{notch} centrado en la frecuencia de la red eléctrica, típicamente 60 Hz. Su función es eliminar una interferencia muy específica y persistente que puede superponerse a la señal fisiológica. Esta clase de filtro resulta útil cuando el entorno de medición introduce ruido periódico por acoplamiento eléctrico. En conjunto, ambas etapas permiten obtener una señal más estable y con menor contaminación, sin destruir la forma global del ECG.

En forma discreta, si $x[n]$ es la muestra cruda de entrada, la cadena de procesamiento puede representarse como
\[
x_{pb}[n] = \mathcal{F}_{pb}\{x[n]\}, \qquad x_f[n] = \mathcal{F}_{notch}\{x_{pb}[n]\},
\]
donde $x_{pb}[n]$ es la salida del pasa-banda y $x_f[n]$ la señal final usada para detección de picos.

Cada etapa causal puede modelarse con un filtro IIR de la forma general
\[
y[n] = \sum_{i=0}^{M} b_i\,x[n-i] - \sum_{j=1}^{N} a_j\,y[n-j],
\]
con función de transferencia
\[
H(z)=\frac{\sum_{i=0}^{M} b_i z^{-i}}{1+\sum_{j=1}^{N} a_j z^{-j}}.
\]
Esta expresión deja explícito que la salida depende solo de la muestra actual y de muestras pasadas, cumpliendo causalidad.

Para el pasa-banda (0.5--40 Hz), una realización típica de segundo orden se expresa como
\[
H_{pb}(z)=\frac{b_0+b_1 z^{-1}+b_2 z^{-2}}{1+a_1 z^{-1}+a_2 z^{-2}},
\]
donde los coeficientes se calculan a partir de la frecuencia de muestreo $f_s$ y de los cortes $f_L=0.5$ Hz y $f_H=40$ Hz.

Para el filtro \textit{notch} en $f_0=60$ Hz, una forma estándar es
\[
H_{notch}(z)=\frac{1-2\cos(\omega_0)z^{-1}+z^{-2}}{1-2r\cos(\omega_0)z^{-1}+r^2 z^{-2}},
\qquad \omega_0=\frac{2\pi f_0}{f_s},
\]
con $0<r<1$ para controlar el ancho de banda de rechazo. Su ecuación en diferencias es
\[
y[n]=x[n]-2\cos(\omega_0)x[n-1]+x[n-2]+2r\cos(\omega_0)y[n-1]-r^2 y[n-2].
\]

En la implementación de este proyecto, el estado interno de cada filtro (muestras previas de entrada y salida) se conserva entre bloques. Esta persistencia evita discontinuidades al procesar tramas sucesivas y garantiza continuidad temporal en la señal filtrada.

\subsection{Visualización de las tres etapas de la señal}

Una característica importante del sistema es que no presenta únicamente la señal final filtrada, sino que conserva y muestra tres representaciones simultáneas de la misma adquisición:

\begin{itemize}
    \item la señal cruda proveniente del microcontrolador,
    \item la señal luego del filtro pasa-banda,
    \item la señal final después del filtro \textit{notch}, sobre la que se realiza la detección de picos.
\end{itemize}

Esta decisión tiene un valor técnico y pedagógico. Desde el punto de vista técnico, permite verificar en qué punto del procesamiento se mejora la señal y si cada filtro está cumpliendo su función esperada. Desde el punto de vista didáctico, facilita comparar visualmente la señal original con sus versiones tratadas, lo que hace más claro el efecto del preprocesamiento digital. Además, cuando aparece un error de adquisición o una interferencia fuerte, la comparación entre curvas ayuda a identificar si el problema proviene del sensor, del canal de transmisión o del propio filtrado.

La representación gráfica se organiza mediante una ventana temporal deslizante. Esto significa que solo se muestra una porción reciente de la señal, mientras las muestras más antiguas salen de la visualización a medida que llegan nuevas. Este mecanismo es el comportamiento típico de un monitor fisiológico, ya que permite seguir el ritmo cardiaco en tiempo real y observar la morfología de la señal sin saturar la pantalla con información acumulada.

\subsection{Detección de picos R}

La detección de picos R se realiza sobre la señal final filtrada, porque en esa etapa el complejo QRS aparece más definido y con menor nivel de ruido. El criterio general consiste en localizar máximos locales que superen un umbral adaptativo. En este enfoque, el umbral no es fijo, sino que depende del comportamiento estadístico de la señal en la ventana más reciente. Esto es importante porque la amplitud del ECG puede variar entre pacientes, entre derivaciones y a lo largo del tiempo.

Si $x[n]$ representa la señal filtrada, el umbral puede expresarse de forma general como
\[
T = \mu_x + k\sigma_x,
\]
donde $\mu_x$ es la media de la señal en la ventana analizada, $\sigma_x$ es su desviación estándar y $k$ es un factor de sensibilidad. Un valor de $k$ más alto exige picos más pronunciados, mientras que un valor menor vuelve la detección más sensible. Esta estrategia permite adaptar el detector a cambios de amplitud y nivel de ruido sin recalibrar manualmente cada caso.

Una vez calculado el umbral, se identifican los máximos locales que lo superan y se aplica un período refractario. El período refractario impide contar dos veces un mismo latido, ya que el complejo QRS puede generar más de una oscilación cercana o puede verse acompañado por artefactos transitorios. Así, solo se acepta un pico nuevo si ha transcurrido el tiempo mínimo establecido desde la última detección válida. El resultado final es una secuencia temporal de instantes asociados a los picos R detectados.

Formalmente, un índice $n_r$ se acepta como pico R si cumple
\[
x_f[n_r-1] < x_f[n_r] \ge x_f[n_r+1], \qquad x_f[n_r] > T,
\]
y además respeta la condición refractaria
\[
n_r-n_{r,\text{prev}} > N_{ref}, \qquad N_{ref}=\left\lfloor t_{ref}\,f_s\right\rfloor,
\]
donde $t_{ref}$ es el tiempo refractario mínimo (por ejemplo, 200 ms) y $f_s$ la frecuencia de muestreo.

\subsection{Estimación de la frecuencia cardiaca}

La frecuencia cardiaca se obtiene a partir de los intervalos RR, es decir, de la diferencia temporal entre dos picos R consecutivos. Si $RR_i$ es el intervalo entre latidos, la frecuencia cardiaca en latidos por minuto puede estimarse como
\[
\mathrm{BPM} = \frac{60}{\overline{RR}},
\]
donde $\overline{RR}$ representa una medida robusta del conjunto reciente de intervalos, por ejemplo la mediana. El uso de la mediana es conveniente porque reduce la influencia de valores atípicos y hace la estimación más estable frente a falsos positivos o pequeñas irregularidades en la detección.

En tiempo discreto, si los picos detectados ocurren en índices $n_{r,i}$, el intervalo puede obtenerse como
\[
RR_i=\frac{n_{r,i}-n_{r,i-1}}{f_s},
\]
y la estimación robusta empleada en línea queda
\[
\widehat{\mathrm{BPM}}[k]=\frac{60}{\mathrm{mediana}\left( RR_{k-m+1},\ldots,RR_k\right)}.
\]

Este cálculo se actualiza en cada bloque procesado, por lo que la visualización no solo muestra la señal ECG, sino también un estimado en tiempo real de la actividad cardiaca. En una presentación o práctica de laboratorio, este valor resulta especialmente útil porque resume de manera inmediata el estado del ritmo y permite observar cómo cambia conforme se modifica la señal adquirida.

\subsection{Actualización de la interfaz gráfica}

La visualización se implementa con un contenedor gráfico que organiza las tres curvas en una misma ventana. El eje horizontal se sincroniza para que las tres gráficas se desplacen de forma coordinada, manteniendo la misma referencia temporal. De esta manera, el usuario puede comparar al mismo instante la señal cruda, la señal pasa-banda y la señal final filtrada con picos R marcados.

La interfaz se actualiza de forma periódica a alta frecuencia para dar continuidad visual. En términos prácticos, esto obliga a equilibrar la fluidez del repintado con el costo computacional del redibujado. Por esa razón, el sistema trabaja sobre ventanas acotadas de muestras recientes y reutiliza los datos acumulados en memoria para evitar recalcular toda la señal en cada refresco.

\subsection{Registro de captura y reutilización de datos}

Además de mostrar la señal en vivo, el sistema puede almacenar la adquisición en un archivo CSV. El registro guarda información de cada muestra junto con la salida filtrada y la estimación de BPM. Este almacenamiento permite posteriormente revisar la señal, analizar su calidad, comparar sesiones o reutilizar los datos como entrada para otros módulos del proyecto.

El soporte de captura y simulación convierte al sistema en una herramienta flexible: puede operar conectado al hardware o reproducir una señal ya registrada, manteniendo el mismo flujo de procesamiento. Esto es útil tanto para validación técnica como para exposición académica, ya que permite demostrar el comportamiento del sistema incluso si no se dispone del microcontrolador en ese momento.

\subsection{Síntesis del enfoque}

En síntesis, el marco conceptual del sistema integra adquisición serie, filtrado causal en dos etapas, detección adaptativa de picos R, estimación de frecuencia cardiaca y visualización simultánea de tres niveles de procesamiento. Esta arquitectura modular facilita comprender cómo se transforma la señal ECG desde su forma original hasta su versión lista para análisis clínico básico. Además, constituye una base adecuada para futuras extensiones, como extracción de características, clasificación automática o integración con modelos de aprendizaje automático.

% --- 8. RESULTADOS Y ANÁLISIS ---
\section{Análisis y Resultados}

\subsection{Enfoque de evaluación del sistema}

El resultado principal de este proyecto no se limita a la construcción de un
clasificador, sino al desarrollo de un sistema funcional de monitoreo ECG en
tiempo real. En consecuencia, la evaluación se centró en verificar cuatro
capacidades del módulo implementado en \texttt{dcp.py}: (i) adquisición continua
desde hardware o desde archivos de simulación, (ii) filtrado causal apto para
procesamiento en línea, (iii) detección estable de picos R con estimación de BPM,
y (iv) visualización simultánea de las distintas etapas de la señal.

Desde esta perspectiva, el modelo de predicción se considera una extensión del
sistema, útil para explorar aplicaciones futuras de análisis automático, pero no
como el eje central del proyecto. Por ello, los resultados se presentan primero
desde el desempeño del monitor en vivo y después, de forma complementaria, desde
la etapa de clasificación.

\subsection{Resultados del monitoreo ECG en tiempo real}

El sistema desarrollado permitió adquirir y procesar la señal ECG de manera
continua a una frecuencia nominal de muestreo de aproximadamente $f_s = 360$~Hz.
La arquitectura implementada admitió dos modos de operación equivalentes desde el
punto de vista del procesamiento:

\begin{itemize}
    \item \textbf{Modo en vivo:} recepción de muestras por puerto serie desde el
        microcontrolador.
    \item \textbf{Modo simulación:} reproducción temporal de archivos CSV ya
        capturados, conservando el mismo flujo interno de filtrado,
        detección y visualización.
\end{itemize}

Esta dualidad constituye un resultado importante, porque hace al sistema robusto
para demostración, depuración y validación experimental. En particular, permite
evaluar el comportamiento del monitor aun cuando no se dispone del hardware en el
instante de la prueba.

\subsection{Desempeño del preprocesamiento causal}

Uno de los aportes más relevantes del sistema es el uso de un pipeline de
preprocesamiento completamente compatible con tiempo real. A diferencia del
procesamiento \textit{offline}, en el cual puede emplearse la señal completa, en
este proyecto el filtrado se ejecuta sobre bloques sucesivos manteniendo el
estado interno de los filtros entre actualizaciones.

La cadena implementada consta de dos etapas:

\begin{enumerate}
    \item filtro pasa-banda Butterworth entre 0.5 y 40~Hz,
    \item filtro \textit{notch} centrado en 60~Hz para atenuar interferencia de red.
\end{enumerate}

En las pruebas realizadas, esta estrategia produjo una mejora visual clara de la
señal respecto a la adquisición cruda. La primera etapa redujo la deriva de línea
base y el ruido de alta frecuencia, mientras que la segunda disminuyó el contenido
periódico asociado al entorno eléctrico. Como resultado, la señal final presentó
complejos QRS más definidos y más adecuados para la detección automática de
latidos.

Un aspecto destacable es que el sistema no solo calcula la señal filtrada final,
sino que conserva en pantalla tres representaciones sincronizadas: señal cruda,
salida del pasa-banda y salida posterior al \textit{notch}. Esto permitió validar
experimentalmente el efecto individual de cada etapa del procesamiento y facilitó
la interpretación del comportamiento del monitor durante las pruebas.

\subsection{Detección de picos R y estimación de frecuencia cardiaca}

La detección de picos R se realizó sobre la señal final filtrada mediante un
umbral adaptativo basado en estadísticos locales de la ventana reciente,
siguiendo el criterio general
\[
T = \mu + 1.4\,\sigma,
\]
combinado con un período refractario de $0.3$~s para evitar dobles conteos del
mismo complejo QRS.

En la implementación, el umbral se actualiza continuamente a partir de una
historia reciente de aproximadamente 2.5~s, lo que permitió adaptar la detección
a variaciones moderadas de amplitud y ruido sin necesidad de recalibración
manual. A partir de los picos detectados se calcularon los intervalos RR y la
frecuencia cardiaca instantánea mediante la mediana de los últimos intervalos
válidos. Además, si transcurren más de 3~s sin una detección confiable, el valor
de BPM se invalida temporalmente para evitar que la interfaz muestre una medida
obsoleta.

Desde el punto de vista funcional, este comportamiento es relevante porque muestra
que el sistema no solo visualiza la señal, sino que también realiza análisis
fisiológico básico en línea. En la práctica, el monitor logró seguir el ritmo de
la señal reproducida y mantener una estimación coherente del BPM durante la
captura, lo cual valida el flujo adquisición--filtrado--detección--visualización
propuesto en el proyecto.

\subsection{Visualización e interacción del sistema}

La interfaz gráfica constituyó otro resultado central. El monitor organiza la
información en una ventana con tres gráficas sincronizadas y un indicador numérico
de BPM de gran tamaño, actualizado de forma periódica. Esta disposición no cumple
solo una función estética: permite inspeccionar la señal en diferentes niveles de
procesamiento, verificar si los picos R se ubican en posiciones correctas y seguir
en tiempo real el comportamiento del ritmo cardiaco.

La ventana visible de aproximadamente 6~s reproduce una lógica de barrido similar
a la de un monitor fisiológico convencional. Esta decisión mejoró la legibilidad
de la señal y evitó la saturación visual que produciría representar toda la
historia adquirida de manera acumulativa. En conjunto, la interfaz convirtió al
sistema en una herramienta no solo de procesamiento, sino también de observación
e interpretación en tiempo real.

\subsection{Registro de datos y reproducibilidad}

Durante la visualización, el sistema puede almacenar automáticamente cada sesión
en archivos CSV con marca temporal. En estos archivos se registra el índice de
muestra, el tiempo, la señal cruda, la señal filtrada y la estimación de BPM.
Este mecanismo aportó dos ventajas prácticas:

\begin{itemize}
    \item \textbf{Trazabilidad experimental:} cada prueba puede revisarse y
        compararse posteriormente.
    \item \textbf{Reutilización de datos:} las capturas generadas en vivo pueden
        emplearse después para análisis \textit{offline}, elaboración de figuras
        o alimentación de módulos complementarios.
\end{itemize}

Desde el punto de vista metodológico, esta capacidad fortalece la reproducibilidad
del proyecto, porque enlaza el sistema de monitoreo en tiempo real con etapas de
análisis posteriores sin romper la continuidad del flujo experimental.

\subsection{Análisis complementario mediante clasificación automática}

Como extensión del sistema principal, se desarrolló además un módulo de análisis
automático de latidos basado en la base pública MIT-BIH Arrhythmia Database
\cite{mitbih}. Esta etapa tuvo como propósito explorar el potencial del proyecto
para incorporar herramientas de apoyo al diagnóstico, pero sin desplazar el foco
del monitor ECG en tiempo real.

En esta extensión se procesaron 16 registros de MIT-BIH y se extrajeron
\textbf{34\,041 latidos}. Para el entrenamiento se conservaron las 7 clases con
mayor representación, obteniéndose \textbf{33\,947 muestras} útiles. De cada
latido se calcularon \textbf{22 características} de tipo temporal, espectral,
wavelet y HRV, sobre las cuales se entrenó un clasificador \textit{Random Forest}.

El modelo alcanzó una exactitud global de \textbf{98.72\,\%} en el conjunto de
prueba, con muy buen desempeño en las clases dominantes N, L y R. No obstante,
las clases minoritarias como A, + y $\sim$ presentaron menor exhaustividad, lo
que confirma que el principal cuello de botella de esta extensión no es el
algoritmo de clasificación en sí, sino el desbalance de clases y la escasez de
ejemplos representativos para eventos poco frecuentes.

Estos resultados son valiosos como prueba de concepto, pero deben interpretarse
como una capa adicional de análisis. En el contexto del proyecto, el modelo no
reemplaza al monitor implementado en \texttt{dcp.py}; más bien demuestra que la
señal adquirida, filtrada y estructurada por el sistema puede servir como base
para desarrollar herramientas futuras de clasificación automática.

\subsection{Discusión}

Los resultados obtenidos permiten afirmar que el aporte principal del proyecto es
la integración exitosa de adquisición, procesamiento y visualización de ECG en un
entorno de tiempo real. El sistema construido logra:

\begin{enumerate}
    \item recibir muestras de manera continua desde hardware real o desde una
        fuente simulada,
    \item mejorar la calidad de la señal mediante filtrado causal en dos etapas,
    \item detectar picos R de forma adaptativa y estimar la frecuencia cardiaca en línea,
    \item registrar la información generada para análisis y reutilización posterior.
\end{enumerate}

Desde una perspectiva de ingeniería, este resultado tiene más peso que la sola
obtención de una alta exactitud en clasificación, porque resuelve el problema
central del proyecto: convertir una adquisición biológica ruidosa en una señal
interpretable y útil en tiempo real. El clasificador, aunque prometedor, depende
de que esta primera capa funcione correctamente.

Las principales limitaciones observadas fueron las siguientes:

\begin{enumerate}
    \item \textbf{Dependencia de la calidad de adquisición:} artefactos por
        movimiento, contacto deficiente de electrodos o ruido del entorno pueden
        degradar la detección si la señal de entrada se deteriora significativamente.
    \item \textbf{Validación clínica acotada:} las capturas propias utilizadas en
        las pruebas corresponden principalmente a condiciones normales, por lo que
        la validación en vivo de arritmias reales sigue siendo limitada.
    \item \textbf{Desbalance en la extensión de aprendizaje automático:} las clases
        minoritarias de MIT-BIH continúan siendo el principal reto para una
        clasificación más uniforme.
\end{enumerate}

En síntesis, el proyecto queda sólidamente sustentado como un sistema de monitoreo
ECG en tiempo real con capacidad de preprocesamiento, análisis básico y registro
de señal. Sobre esa base, la clasificación automática aparece de manera natural
como una línea de expansión futura, y no como el núcleo de la solución.

% =============================================================
%  SECCIÓN: CONCLUSIONES
% =============================================================

\section{Conclusiones}

El presente proyecto permitió desarrollar y validar un sistema funcional de
monitoreo ECG en tiempo real, capaz de integrar adquisición de señal,
preprocesamiento digital, detección de picos R, estimación de frecuencia
cardiaca y visualización simultánea de las distintas etapas del procesamiento.
Desde esta perspectiva, el principal aporte del trabajo no radica únicamente en
la obtención de buenas métricas de clasificación, sino en haber construido una
cadena completa que transforma una señal biológica ruidosa en información útil,
interpretable y reutilizable en línea.

En el plano experimental, el sistema demostró un comportamiento consistente tanto
en modo de adquisición en vivo como en modo de simulación a partir de archivos
CSV. La arquitectura implementada permitió trabajar con una frecuencia de
muestreo cercana a 360~Hz, mantener la continuidad del filtrado causal entre
bloques y obtener una estimación establ e del BPM a partir de los intervalos RR.
Asimismo, la posibilidad de visualizar en paralelo la señal cruda, la salida del
filtro pasa-banda y la señal posterior al filtro \textit{notch} aportó valor
técnico y didáctico, al facilitar la validación del procesamiento y la
interpretación del comportamiento del sistema durante las pruebas.

Como extensión del sistema principal, el módulo de clasificación automática de
latidos mostró resultados altamente competitivos. El clasificador
\textit{Random Forest}, entrenado sobre 16 registros de la base MIT-BIH a partir
de 22 características temporales, espectrales, wavelet y HRV, alcanzó una
exactitud global de \textbf{98.72\,\%} en el conjunto de prueba. Este resultado
confirma que la señal procesada y estructurada por el sistema constituye una base
adecuada para tareas posteriores de análisis automático. Sin embargo, el menor
desempeño en clases minoritarias como A, + y $\sim$ evidencia que el principal
desafío pendiente sigue siendo el desbalance de clases y la limitada
representatividad de eventos poco frecuentes.

Las limitaciones identificadas también permiten delimitar con claridad el alcance
actual del proyecto. Por un lado, la calidad del monitoreo continúa dependiendo
de condiciones físicas de adquisición, como el contacto de electrodos, el ruido
eléctrico del entorno y los artefactos por movimiento. Por otro, la validación
en vivo se realizó principalmente sobre capturas compatibles con ritmo normal,
por lo que todavía no es posible afirmar un desempeño clínico robusto frente a
arritmias reales adquiridas en tiempo real. En consecuencia, el sistema debe
entenderse como una plataforma sólida de desarrollo y prueba, más que como una
herramienta diagnóstica concluida.

Como líneas de trabajo futuro, resulta razonable priorizar la mejora del
hardware de adquisición para reducir artefactos, ampliar el conjunto de capturas
propias con señales patológicas reales o clínicamente validadas, fortalecer el
tratamiento de clases minoritarias dentro del módulo de aprendizaje automático y
explorar arquitecturas más avanzadas, como redes convolucionales de una dimensión
o modelos recurrentes. Del mismo modo, una evaluación comparativa con criterio
experto permitiría complementar las métricas estadísticas con indicadores de
relevancia clínica.

En conjunto, el trabajo establece una base reproducible y técnicamente coherente
para el desarrollo de sistemas de monitoreo cardíaco inteligente de bajo costo.
Su mayor fortaleza reside en haber unido instrumentación, procesamiento digital y
análisis computacional dentro de una misma plataforma, dejando abierto un camino
natural hacia aplicaciones futuras en telemedicina, monitoreo portátil y apoyo
al análisis clínico asistido por computadora.
